\chapter{Introduction and Motivation}

With the continued development of embodied intelligence and humanoid robotics, research has gradually moved from locomotion in structured environments toward operation in complex, dynamic, and uncertain settings. Among the many challenging benchmarks in this direction, humanoid parkour stands out as a representative task because it requires the robot to step over obstacles, climb stairs, traverse gaps, absorb impacts, and recover balance in rapid succession. Compared with standard flat-ground walking, parkour places much stronger demands on perception, planning, control, and whole-body coordination. A robot must not only remain stable, but also maintain motion continuity under changing terrain and respond in real time to disturbances. For this reason, humanoid parkour has become an important testbed for evaluating the agility, robustness, and environmental adaptability of humanoid robots.

From the perspective of control, humanoid parkour is not a single locomotion problem, but a long-horizon decision-making problem that couples sensing, temporal reasoning, and motion generation. During traversal, the robot must continuously interpret partial terrain information, assess the feasibility of upcoming actions, and adjust posture and contact strategy accordingly. The state at one moment influences the range of feasible states in the future, which means that successful control requires an explicit understanding of temporal dependencies rather than only instantaneous observations. In real-world environments, however, observations are often incomplete, noisy, and time-varying. Terrain geometry may change, targets may shift, and unexpected disturbances may occur. These factors make it difficult for conventional methods based on fixed trajectories, short-term optimization, or manually designed rules to achieve reliable performance in complex parkour scenarios.

The difficulty is further amplified in non-structured environments with uncertain friction, irregular footholds, dynamic obstacles, and imperfect model parameters. Under such conditions, the robot must do more than remain upright: it must jump accurately, land safely, transition smoothly between motion primitives, and maintain task progress toward a desired goal. This requires a policy that can identify which historical information is still relevant, attend to the most informative environmental cues, and generate coherent actions under goal constraints. Existing control approaches often struggle to balance stability, flexibility, and generalization when the task involves multi-stage behaviors and frequent contact changes. As a result, a more expressive and temporally aware framework is needed to support robust parkour execution.

Motivated by these challenges, this thesis investigates target-driven humanoid robot parkour and explores the integration of goal conditioning, memory augmentation, and attention mechanisms into the control framework. The target-driven formulation encourages the robot to organize intermediate behaviors around a high-level objective, rather than optimizing only for short-term stability. Memory augmentation helps preserve key information across time, enabling the policy to exploit historical context in long-horizon tasks and reducing the risk of forgetting task-relevant state transitions. Attention mechanisms further improve the model's ability to focus on critical motion cues and salient environmental features, allowing it to distinguish useful information from background noise in rapidly changing scenes. By combining these components, the proposed approach aims to produce motions that are more continuous, adaptive, and robust across diverse parkour conditions.

Beyond methodological contributions, this work also has clear practical value. In applications such as disaster response, infrastructure inspection, exploration of confined spaces, and navigation across unstructured terrain, humanoid robots may encounter stairs, slopes, gaps, rubble, and irregular landing surfaces. In such settings, planar walking alone is insufficient. A humanoid robot with parkour capability can make faster decisions, traverse obstacles more effectively, and reach target locations with greater autonomy. Therefore, research on target-driven humanoid parkour not only advances dynamic control and motion generalization, but also contributes to the broader development of embodied intelligence and real-world robotic deployment.

This thesis focuses on the design and evaluation of a target-driven, memory-attention-enhanced control framework for humanoid parkour. The central objective is to investigate how a robot can achieve stable and efficient traversal under complex terrain conditions and dynamic goal constraints. By addressing this problem, the thesis seeks to provide a new perspective for high-agility humanoid locomotion and to support the future deployment of humanoid robots in practical, safety-critical environments.